\SourceLeaf{193}
\Needspace{17\baselineskip}
Beispiel: Nr. 138 Conradi.\quad \(M=65.02,\quad m=46.08\)
\begin{center}
\fontsize{9}{14}\selectfont\setlength{\tabcolsep}{5pt}
\begin{tabular}{lrr}
&4 Proc.&\(3\tfrac{1}{2}\) Proc.\\
\(\log\varphi m\ .\ .\ .\ .\)&1.12875&1.15233\\
\(\tfrac{7}{24}\log\rho\ .\ .\ .\ .\)&\(-497\)&\(-436\)\\
&13.2979&14.0594\\
\(\tfrac{1}{24}\rho^{\frac{5}{4}}\ .\ .\ .\ .\)&0.0397&399\\\cline{2-3}
\(R=\)&13.3376&14.0993\\\cline{2-3}
für \(Q,\ \log\psi\ .\ .\ .\)&0.80912&0.82183\\
\(\log gG\ .\ .\)&\(+897\)&\(+897\)\\
\(\tfrac{1}{4}\log\rho\ .\ .\ .\)&\(-426\)&\(-374\)\\
\(Q=\)&6.5137&6.7150\\
\(R-Q=\)&6.8239&7.3843\\\hline
für \(S,\ \log\psi\ .\ .\ .\)&0.81325&0.82606\\
\(\log gG\ .\ .\)&\(+1744\)&\(+1744\)\\
\(S=\)&6.7716&6.9743\\\hline
\end{tabular}
\end{center}
\[
\begin{aligned}
\log\rho&=-0.0170333393\quad\text{für 4 Proc.}\\
\log\rho&=-0.0149403498\quad\text{für }3\tfrac{1}{2}\text{ Proc.}
\end{aligned}
\]
\begin{center}
\fontsize{9}{14}\selectfont\setlength{\tabcolsep}{4pt}
\begin{tabular}{c|c}
für \(Q\)&für \(S\)\\
\(\psi(46.04;64.65)\ .\ .\ 18.61\)&\(\psi(45.08;64.52)\ .\ .\ 19.44\)\\
18.61&19.44\\
\begin{tabular}{r||r|r}64&0.82079&0.83374\\65&0.80284&0.81540\end{tabular}&\begin{tabular}{r|r}0.82251&0.83552\\0.80470&0.81732\end{tabular}\\
\(g(46.56)\ .\ .\ .\ 25\)&\(g(46.08)\ .\ .\ .\ 588\)\\
\(G(65.33)\ .\ .\ .\ 872\)&\(G(65.27)\ .\ .\ .\ 1156\)
\end{tabular}
\end{center}

\noindent\hspace*{1.2em} Bei der Aufstellung dieser Tafeln hat \textsc{Gauss} sich seiner fünfstelligen Tafel zur Berechnung des Logarithmus der Summe von Grössen, die nicht selbst, sondern nur durch ihre Logarithmen gegeben sind, bedient und deshalb bei jedem Zinsfuss und jedem Altersunterschiede mit der Bestimmung der Rentenwerthe für die höchsten Lebensjahre den Anfang gemacht. Die einzelnen Zahlenangaben können aus diesem Grunde auch abgesehen von der Verbesserung, welche die Sterblichkeitstafel durch erweiterte Erfahrungen der Preussischen Witwenkasse schon seither erlitten hat, um einzelne Einheiten in der fünften Decimale der Logarithmen ungenau sein.

Einige Rechenfehler, auf die ich durch Bildung der Quotienten zwischen den \(3\tfrac{1}{2}\) und 4 procentigen Rentenwerthen und der Differenzen der auf einander folgenden Quotienten aufmerksam geworden bin, habe ich beim Abdruck berichtigt. In Einheiten der fünften Decimale des Logarithmus betragen diese Fehler an den Orten ihres Entstehens: \(+20\) für das 77. Jahr der Frau in deren Leibrentenwerthen bei 4 Proc. ferner für das 71. und 93. Jahr des Mannes und die resp. Altersunterschiede der Frau von \(+1\) und 0 Jahr bei \(3\tfrac{1}{2}\) Proc.; \(+10\) für das 76. Jahr des Mannes in dessen Leibrentenwerth bei \(3\tfrac{1}{2}\) Proc. und ebenso viel für das 79. Jahr des Mannes und den Altersunterschied der Frau von \(-9\) Jahr bei 4 Proc.; \(-10\) für das 28. 68. 91. 94. und 94. Jahr des Mannes und die resp. Altersunterschiede von \(-8\), \(-13\), \(-11\), \(-2\), und \(-14\) Jahr bei den resp. Procenten \(4.4.3\tfrac{1}{2}.4\) und \(3\tfrac{1}{2}\). Diese und einige andere Rechenfehler von geringerem Betrage haben auf die Bestimmung der Rentenwerthe für die zunächst jüngeren Altersjahre einigen Einfluss gehabt, der allmählig und im äussersten Falle erst für das Ende des zweiten Jahrzehnt verschwindet. Die Angaben der einfachen Leibrenten in Zahlen sind aus den Logarithmen abgeleitet und haben nach der angegebenen Berichtigung derselben hier auch eine entsprechende Abänderung erfahren müssen.

Die bei den Anwendungen der Tafeln zu gebrauchenden Formeln und die Rechnungsbeispiele zu denselben sind den zerstreuten Notizen auf einzelnen Handblättchen entlehnt und hier durch einige Einschaltungen erläutert
\begin{flushright}\textsc{Schering.}\end{flushright}
\begin{center}\rule{.3\linewidth}{.3pt}\end{center}
